\cleardoublepage
\chapter{Supplementary Materials}

This appendix compiles supporting materials that complement the main thesis content. The items included here are intended to improve reproducibility and provide implementation details that are too lengthy for the core chapters.

\section{System Configuration Summary}
Table~\ref{tab:appendix-system-config} summarizes the software and hardware setup used during development and evaluation.

\begin{table}[h]
    \centering
    \caption{System configuration used in this study}
    \label{tab:appendix-system-config}
    \begin{tabular}{ll}
        \hline
        Item & Specification \\
        \hline
        Operating System & [Specify operating system] \\
        CPU & [Specify processor model] \\
        RAM & [Specify memory size] \\
        GPU & [Specify GPU model or N/A] \\
        Programming Language & [Specify language and version] \\
        Key Libraries/Frameworks & [Specify libraries and versions] \\
        \hline
    \end{tabular}
\end{table}

\section{Additional Experimental Settings}
The following settings were used to ensure fair and repeatable experiments:

\begin{itemize}
    \item Dataset split ratio: [e.g., 70\% training, 15\% validation, 15\% testing].
    \item Random seed: [Specify seed value].
    \item Number of runs per experiment: [Specify number].
    \item Evaluation metrics: [Specify metrics, e.g., accuracy, precision, recall, F1-score].
\end{itemize}

\section{Extended Results}
Table~\ref{tab:appendix-extended-results} provides additional performance results that support the analysis in Chapter~4.

\begin{table}[h]
    \centering
    \caption{Extended performance results}
    \label{tab:appendix-extended-results}
    \begin{tabular}{lccc}
        \hline
        Method & Metric 1 & Metric 2 & Metric 3 \\
        \hline
        Baseline & [Value] & [Value] & [Value] \\
        Proposed Method & [Value] & [Value] & [Value] \\
        \hline
    \end{tabular}
\end{table}

\section{Implementation Notes}
Key implementation details are listed below for completeness:

\begin{enumerate}
    \item Data preprocessing steps: [Describe main preprocessing pipeline].
    \item Model training procedure: [Describe optimizer, learning rate, epochs, batch size].
    \item Inference workflow: [Describe inference pipeline and post-processing].
    \item Error handling and edge cases: [Describe notable cases and handling strategy].
\end{enumerate}

\section{Usage Instructions for Reproduction}
To reproduce the reported results, follow these steps:

\begin{enumerate}
    \item Prepare the environment by installing the required dependencies.
    \item Download or generate the dataset according to the project instructions.
    \item Run training with the same hyperparameters listed in this thesis.
    \item Evaluate the trained model on the test set.
    \item Compare results with the values reported in Chapter~4 and this appendix.
\end{enumerate}

\clearpage